\documentclass[10pt,a4paper]{amsart}
\usepackage[margin=19mm]{geometry}
\usepackage{amsmath,amssymb,mathtools}
\usepackage[scr=rsfs,cal=euler]{mathalfa}
\usepackage{xfrac}
\usepackage[unicode,hidelinks,bookmarks=false]{hyperref}
\newcommand{\ie}{i.e.\ }
\newcommand{\cf}{cf.\ }
\newcommand{\resp}{respectively\ }
\newcommand{\Subsection}{\subsection}
\providecommand{\nobreakdash}{\nobreak\hbox{-}}
\setlength{\emergencystretch}{2em}
\multlinegap0pt
\usepackage{float}
\usepackage{xcolor}
\usepackage{amssymb}
\usepackage{bm}
\usepackage{dsfont}
\usepackage{tikz}
\usepackage{siunitx}
\usepackage[shortlabels]{enumitem}
\newtheorem{proposition}{Proposition}
\title{Arbitrary order stationarity preserving stabilized finite elements for multidimensional nonlinear hyperbolic problems. \\ Application to the Euler equations with gravity}
\author{Moussa Ziggaf, Davide Torlo, Mario Ricchiuto}
\date{Source: arXiv:2603.23185v1 (CC BY 4.0)}
\begin{document}
\maketitle

\noindent\emph{Source and licence.} Arbitrary order stationarity preserving stabilized finite elements for multidimensional nonlinear hyperbolic problems. \\ Application to the Euler equations with gravity. Version arXiv:2603.23185v1 (CC BY 4.0), distributed by arXiv under the Creative Commons Attribution 4.0 International licence, http://creativecommons.org/licenses/by/4.0/, as recorded in the arXiv licence metadata for this exact version and shown on the abstract page https://arxiv.org/abs/2603.23185v1. Full licence notice: \texttt{licenses/arxiv-2603.23185-CC-BY-4.0.txt}.\par\smallskip

\noindent\textit{Editorial scope.} Counted unit: the article's proposition prop:equilibria\_local (source0.tex lines 1632--1636). The article's complete original proof of it is delivered with it (source0.tex lines 1637--1642, one \texttt{proof} environment of 6 lines). No mathematics is written, repaired or re-derived: the statement and the proof are the article's own lines, in the article's own notation, and no retained line is substituted or re-flowed.  Retained uncounted as the source-local context the proof needs: lines 945--949; lines 1453--1465; lines 1507--1512; lines 1604--1614; lines 1623--1629.  Nothing was omitted from any retained span, so the package carries no omission record.  No retained line contains a citation, so the wrapper carries no bibliography and no bibliography entry.  No retained line contains a figure environment, an image-inclusion command or drawing code, so no image is delivered and the package declares no asset dependency.  Editorial formatting only, in addition: an AMS \texttt{amsart} preamble that restates the article's own declarations and macros used by the retained lines, namely the environments \texttt{proposition} on separate counters, together with the standard packages those lines need.  The retained environments print in this document as Proposition 1 in order of appearance, whereas the article prints the same items with the numbering of its own full document.  The complete native source of the article as distributed in the arXiv e-print for this exact version is bundled unchanged as \texttt{sources/197060/source0.tex}.
\begin{equation}\label{eq:weakG1-SU}
M \dfrac{d }{dt} \mathrm{W} +
 \textsf{div}_h \mathbf{F}    + \mathsf{st}_h(\mathrm{F}_1,\mathrm{F}_2)=
M  \mathrm{S},
\end{equation}

\begin{equation}
	\label{eq:DeC-GF}
	\begin{split}
	     \mathrm{W}^{m,(k)}  =            \mathrm{W}^n
	- &\Delta t \sum_{r=0}^M \theta_r^m \mathsf{M}^{-1}  M^{\textsf{SU}}(\mathrm{W}^{r,(k-1)}) \dfrac{d\mathrm{W^{(k-1)}}(t^r)}{dt} \\
	-& \Delta t \sum_{r=0}^M \theta_r^m \mathsf{M}^{-1}  \{ \textsf{R}_h(\mathrm{W}^{r,(k-1)})    + \mathsf{ST}_h(\mathrm{W}^{r,(k-1)}) \}
%	     \mathcal{L}^{2,m}( \mathsf{W}^{(k-1)} )%\\
%	    \mathcal{L}^{2,m}( \mathsf{W}^{k-1} ) =  & M_1\otimes M_2  (\mathrm{W}^{m,(k-1)} -  \mathrm{W}^0 )
%	+ \Delta t \sum_{r=0}^M \theta_r^m \{ M^{\textsf{SU}}(\mathrm{W}^{r,(k-1)}) \dfrac{d\mathrm{W}(t^r)}{dt} 
% + \textsf{DIV}_h \mathbf{F}(\mathrm{W}^{r,(k-1)})    + \mathsf{ST}_h(\mathrm{W}^{r,(k-1)} )  -  M(\mathrm{W}^{r,(k-1)} )
%	 \mathrm{S}(\mathrm{W}^{r,(k-1)})\}  
	 \end{split}
\end{equation}

 \begin{equation}\label{eq:Phi_gf}
 \begin{aligned}
  \partial_{x^1} \partial_{x^2} \mathcal{F}_{1,h} + \partial_{x^2}\partial_{x^1} \mathcal{F}_{2,h} =   & \; \partial_{x^1} \partial_{x^2}\Psi^{ij}_h\;, \text{ for }\mathbf{x} \in E_{ij},  \\
[  \Psi^{ij}_h]^{ij}_{pq} := &\;[\mathcal{F}_{1,h} - \mathcal{F}_{1,h}(x^1_{i,0},x^2) + \mathcal{F}_{2,h} - \mathcal{F}_{2,h}(x^1, x^2_{j,0})]_{pq} =  \int_{x^2_{j,0}}^{x^2_{j,q}} \int_{x^1_{i,0}}^{x^1_{i,p}} \nabla\cdot\mathbf{F}_h\,d\mathbf{x}
  \end{aligned}
 \end{equation}

\begin{proposition}\label{prop:equilibria}
For any arbitrary  univariate functions $\mathbf{f}(x^1)$ and $\mathbf{g}(x^2)$, with discrete interpolated counterparts $\mathsf{f}$ and $\mathsf{g}$,
the general class of functions $\mathbf{W}_h$ such that their integrated fluxes and source can be written as 
\begin{equation}\label{eq:equilibria}
	\begin{split}
	 \mathcal{F}_{1,h}(\mathbf{W}_h)(\mathbf{x}) + \mathcal{F}_{2,h}(\mathbf{W}_h)(\mathbf{x}) + \mathcal{S}_h(\mathbf{W}_h)(\mathbf{x}) = &\; \mathbf{f}_h(x^1)+\mathbf{g}_h(x^2) , \text { or discretely}\\[5pt]
		(\mathcal{F}_{1})^{ij}_{pk} + (\mathcal{F}_{2})^{ij}_{pk} + \mathcal{S}^{ij}_{pk} =& \; \mathsf{f}^i_p + \mathsf{g}^j_k,
	\end{split}
\end{equation} 
describe discrete stationary solutions of \eqref{eq:weakG1-SU} and \eqref{eq:DeC-GF}.
\end{proposition}

 \begin{equation}\label{eq:equilibria-phys}
 \begin{aligned}
 &\mathbf{g}_h(x^2) =  \mathcal{F}_{1,h}(\mathbf{W}_h)(x^1 =x^1_0,x^2)\;, \quad   \mathbf{f}_h(x^1) =\mathcal{F}_{2,h}(\mathbf{W}_h)(x^1, x^2 =x^2_0)\\
\Rightarrow & \int_{x^2_0}^{x^2}[\mathbf{F}_{1,h}-\mathbf{F}_{1,h}(x^1 =x^1_0,x^2) ] \mathrm{d}x^2 + \int_{x^1_0}^{x^1}[\mathbf{F}_{2,h}-\mathbf{F}_{2,h}(x^1,x^2=x^2_0) ] \mathrm{d}x^1=
 \int_{x^1_0}^{x^1} \int_{x^2_0}^{x^2}\mathbf{S}_hd\mathbf{x}
\end{aligned}
 \end{equation} 

\begin{proposition}\label{prop:equilibria_local}
The class of functions $\mathbf{W}_h$ belonging to the kernel of Proposition \ref{prop:equilibria} 
with $\mathbf f$ and $\mathbf g$ given by \eqref{eq:equilibria-phys}, also verifies    $\Psi_h+\mathcal{S}_h=0$  $\forall E_{ij}\in \Omega$,
with $\Psi_h$  the  array of integrated divergences  \eqref{eq:Phi_gf}.
\end{proposition}

\begin{proof}
 The result is shown iteratively. 
 Within this first element $E_{00}$ we integrate on each subcell $[x^1_{0,0},x^1_{0,p}]\times [x^2_{0,0},x^2_{0,q}]$ with increasing $p$ and $q$ from $0$ to $K$.
 This allows us to obtain $\Psi_h+\mathcal{S}_h=0$  for $E_{00}$. We then  increase $i$ and $j$ by 1 in each direction, and use
 the result on $E_{00}$ to prove the same on $E_{10}$ and $E_{01}$. We continue iterating on $i$ and $j$ until all $E_{ij}$ have been covered.
\end{proof}

\end{document}
